\documentclass[11pt,twocolumn,a4paper]{article}

% ─── Packages ─────────────────────────────────────────────────────────
\usepackage[margin=2.0cm,top=2.4cm,bottom=2.4cm]{geometry}
\usepackage[utf8]{inputenc}
\usepackage[T1]{fontenc}
\usepackage{lmodern}
\usepackage{microtype}
\usepackage{amsmath,amssymb,amsthm,mathtools}
\usepackage{physics}
\usepackage{bm}
\usepackage{graphicx}
\usepackage{booktabs}
\usepackage{array}
\usepackage{multirow}
\usepackage{enumitem}
\usepackage{float}
\usepackage{caption}
\usepackage{algorithm}
\usepackage{algpseudocode}
\usepackage{natbib}
\usepackage{xcolor}
\usepackage{mathrsfs}
\usepackage[colorlinks=true,linkcolor=blue!60!black,
            citecolor=blue!60!black,urlcolor=blue!60!black]{hyperref}
\usepackage{fancyhdr}
\usepackage{titlesec}

% ─── Page Style ───────────────────────────────────────────────────────
\pagestyle{fancy}
\fancyhf{}
\fancyhead[L]{\small\textit{S-Entropy Embedding and the Forced Measurement Protocol}}
\fancyhead[R]{\small\thepage}
\renewcommand{\headrulewidth}{0.4pt}

\titleformat{\section}{\large\bfseries}{\thesection.}{0.5em}{}
\titleformat{\subsection}{\normalsize\bfseries}{\thesubsection.}{0.5em}{}
\titleformat{\subsubsection}{\normalsize\itshape}{\thesubsubsection.}{0.5em}{}

% ─── Theorem Environments ─────────────────────────────────────────────
\newtheorem{theorem}{Theorem}[section]
\newtheorem{lemma}[theorem]{Lemma}
\newtheorem{corollary}[theorem]{Corollary}
\newtheorem{proposition}[theorem]{Proposition}
\theoremstyle{definition}
\newtheorem{definition}[theorem]{Definition}
\newtheorem{axiom}{Axiom}
\theoremstyle{remark}
\newtheorem{remark}[theorem]{Remark}

% ─── Custom Commands ──────────────────────────────────────────────────
\newcommand{\kB}{k_{\mathrm{B}}}
\newcommand{\R}{\mathbb{R}}
\newcommand{\N}{\mathbb{N}}
\newcommand{\Sk}{S_{\mathrm{k}}}
\newcommand{\St}{S_{\mathrm{t}}}
\newcommand{\Se}{S_{\mathrm{e}}}
\newcommand{\Sent}{\bm{S}}
\newcommand{\Ocat}{\mathcal{O}_{\mathrm{cat}}}
\newcommand{\Ophys}{\mathcal{O}_{\mathrm{phys}}}

% ─── Caption Setup ────────────────────────────────────────────────────
\captionsetup{font=small,labelfont=bf}

% ══════════════════════════════════════════════════════════════════════
\begin{document}

\twocolumn[
\begin{@twocolumnfalse}
\begin{center}

{\LARGE\bfseries S-Entropy Embedding Uniqueness and the Forced Measurement Protocol:\\[4pt]
Deriving the State Encoder from Bounded Phase Space Theory}

\vspace{12pt}

{\large Kundai Farai Sachikonye}

\vspace{6pt}

{\small\itshape Chair of Experimental Bioinformatics, Technical University of Munich \\[0.3em]
\small\itshape AIMe Registry for Artificial Intelligence \\[0.2em]
\small\itshape Maximus-von-Imhof-Forum 3, 85354 Freising, Germany \\[0.2em]
\texttt{kundai.sachikonye@wzw.tum.de}}

\vspace{6pt}{\small May 2026}

\vspace{16pt}

\begin{minipage}{0.92\textwidth}
\textbf{Abstract.}
We derive the three-coordinate S-entropy state encoder as a mathematical necessity
of any observer operating on bounded phase space systems, not as an engineering
design choice. The argument proceeds from the Bounded Phase Space axiom through
five theorems to a unique forced measurement protocol. Poincar\'{e} recurrence
forces all persistent bounded systems into oscillatory regimes with discrete
Koopman spectra. Three scalar functionals of this spectrum --- kinematic entropy
$\Sk$, temporal span $\St$, and harmonic fraction $\Se$ --- are shown to be the
minimal sufficient statistics for categorical system identification under the
operations supported by the membrane computing substrate. The map
$\sigma\colon \mathcal{X} \to \lbrack 0,1\rbrack^3$, $\sigma(X) = (\Sk, \St, \Se)$, satisfies
a categorical sufficiency theorem: two systems with identical S-entropy coordinates
are computationally indistinguishable under aperture filtering, cascade routing,
and strobe readout. A timescale separation theorem establishes that $\Sk$, $\St$,
and $\Se$ must be measured on hierarchically ordered timescales
$\tau_k \ll \tau_t \ll \tau_e$, each derived from the system's own spectral
structure. The measurement sequence is therefore forced: any observer that violates
the timescale ordering generates aliased coordinates that fail the sufficiency
condition. The S-entropy space carries a natural product Fisher information metric
$g = g_k \oplus g_t \oplus g_e$ compatible with the partition depth measure from
the bilayer substrate. A conservation law $\sum_k p_k = 1$ holds on the spectral
probability simplex by Parseval's theorem; this law is preserved under the
three-pass measurement protocol. Validation on 39 standard compounds confirms
numerical separability: the six chemical families occupy disjoint regions of
S-entropy space, with inter-family distances 1.09 to 4.38 times larger than
intra-family diameters.
\end{minipage}

\vspace{8pt}
\noindent\textbf{Keywords:}
S-entropy, Koopman spectrum, state encoder, forced measurement, categorical
sufficiency, Fisher information, partition coordinates, Parseval conservation,
timescale separation, bounded phase space

\vspace{18pt}
\end{center}
\end{@twocolumnfalse}
]

%──────────────────────────────────────────────────────────────────────
\section{Introduction}
\label{sec:introduction}
%──────────────────────────────────────────────────────────────────────

Every computing substrate requires a state encoder: a procedure that maps the
physical state of an input system to a representation on which the computer's
operations act. For classical digital systems, the encoder is the
analogue-to-digital converter, whose design is guided by the Nyquist-Shannon
sampling theorem \citep{shannon1949}. For neuromorphic and analogue systems,
the encoder is typically designed by heuristic optimisation against a training
corpus \citep{maass1997}. Neither approach is derived from first principles; both
presuppose a fixed target representation and engineer backward to the encoder.

This paper derives the encoder forward from the physics of bounded phase space.
The starting point is purely mathematical: any persistent physical system occupies
a bounded region of phase space with finite Liouville measure \citep{liouville1838}.
From this single axiom, we show that the system's Koopman spectrum
\citep{koopman1931} is the unique information-complete representation of its state,
and that three scalar functionals of this spectrum --- kinematic entropy $\Sk$,
temporal span $\St$, and harmonic fraction $\Se$ --- are the minimal sufficient
statistics for categorical identification compatible with the membrane computing
operations established in the companion bilayer-substrate paper \citep{companion1}.

The three-coordinate system we derive is not a design choice but a forced
consequence of three constraints operating simultaneously: (i) completeness under
categorical operations, (ii) measurability without demolishing the partition
structure of the substrate, and (iii) computability on the membrane hardware
without external reference signals. We call this the \emph{forced measurement
protocol}: once the Bounded Phase Space axiom is accepted and the membrane
substrate is fixed, every other encoding is either informationally incomplete
or destroys the substrate's partition structure.

The paper is structured as follows. Section~\ref{sec:bps} establishes the
Koopman spectral representation from the Bounded Phase Space axiom.
Section~\ref{sec:sentropy} defines the three S-entropy coordinates and establishes
their basic properties. Section~\ref{sec:sufficiency} proves the categorical
sufficiency theorem. Section~\ref{sec:timescale} establishes the timescale
separation and the non-demolition commutation relation.
Section~\ref{sec:protocol} derives the forced measurement protocol.
Section~\ref{sec:conservation} proves the S-entropy conservation law and the
Fisher metric structure. Section~\ref{sec:implementation} describes the physical
realisation of the three passes in the membrane substrate.
Section~\ref{sec:validation} presents numerical validation.
Sections~\ref{sec:discussion} and~\ref{sec:conclusion} discuss and conclude.

%──────────────────────────────────────────────────────────────────────
\section{Bounded Phase Space and Koopman Spectra}
\label{sec:bps}
%──────────────────────────────────────────────────────────────────────

\subsection{The Foundational Axiom}

\begin{axiom}[Bounded Phase Space]
\label{ax:bps}
Every persistent physical system $\mathcal{X}$ occupies a region
$\Omega \subset \mathbb{R}^{2d}$ of $d$-dimensional phase space satisfying:
\begin{equation}
\mu(\Omega) = \int_\Omega \frac{d^d q\, d^d p}{h^d} < \infty,
\label{eq:bounded}
\end{equation}
where $\mu$ is the Liouville measure normalised by Planck's constant $h$ per
degree of freedom \citep{planck1901, heisenberg1927}. Moreover, $\Omega$ admits
a sequence of refining measurable partitions $\mathscr{P}_0 \supset \mathscr{P}_1
\supset \cdots$ with $|\mathscr{P}_k| \to \infty$ and
$\sup_{P \in \mathscr{P}_k} \mathrm{diam}(P) \to 0$.
\end{axiom}

\noindent Equation~\eqref{eq:bounded} implies that the number of
distinguishable microstates $N_{\max} = \mu(\Omega)$ is finite. This is the
only axiom of the framework; all subsequent results are theorems.

\subsection{Oscillatory Necessity}

\begin{theorem}[Oscillatory Necessity]
\label{thm:osc}
Every system satisfying Axiom~\ref{ax:bps} is generically oscillatory:
its long-term dynamics are exhausted by quasiperiodic trajectories on invariant
tori, with the single exception of a measure-zero set of non-generic initial
conditions.
\end{theorem}

\begin{proof}
Let $\Phi^t: \Omega \to \Omega$ be the Hamiltonian flow on $\Omega$. By Liouville's
theorem, $\Phi^t$ preserves the measure $\mu$. Since $\mu(\Omega) < \infty$, the
Poincar\'{e} Recurrence theorem \citep{poincare1890} applies: for $\mu$-almost
every $\mathbf{x} \in \Omega$, the trajectory $\{\Phi^t(\mathbf{x}): t \geq 0\}$
returns within $\epsilon$ of $\mathbf{x}$ for every $\epsilon > 0$.
We classify the possible long-term behaviours:

\noindent\textbf{Static fixed points}: Form a measure-zero set in a generic
Hamiltonian system (by the implicit function theorem applied to $\nabla H = 0$).

\noindent\textbf{Monotone drift}: If $q_i(t) \to \pm\infty$, this contradicts
$\mu(\Omega) < \infty$. Monotone drift is impossible on a bounded domain.

\noindent\textbf{Non-recurrent chaos}: Positive-measure non-recurrent trajectories
would require infinite orbit volume inside a finite-measure domain: impossible.

\noindent\textbf{Recurrent quasiperiodic}: KAM theory \citep{kolmogorov1954,
arnold1963, moser1962} guarantees that for generic Hamiltonian systems, the
phase space is dominated by invariant tori supporting quasiperiodic motion with
well-defined frequency vectors $\bm{\omega} = (\omega_1, \ldots, \omega_d)$.

After excluding measure-zero exceptions, generic trajectories are quasiperiodic
with well-defined power spectra --- the defining property of oscillatory dynamics.
\end{proof}

\subsection{Discrete Koopman Spectrum}

\begin{theorem}[Koopman Spectral Decomposition]
\label{thm:koopman}
For any system satisfying Axiom~\ref{ax:bps}, the Koopman operator
$(\mathcal{U}^t g)(\mathbf{x}) = g(\Phi^t \mathbf{x})$ on $L^2(\Omega, \mu)$
has a purely discrete spectrum: there exist frequencies $\{\omega_k\}_{k=1}^K$,
amplitudes $\{A_k\}$, and phases $\{\phi_k\}$ such that for any
$f \in L^2(\Omega, \mu)$:
\begin{equation}
f(\Phi^t \mathbf{x}) = \sum_{k=1}^K A_k \cos(\omega_k t + \phi_k) + \varepsilon(t),
\label{eq:koopman}
\end{equation}
with $\|\varepsilon\|_{L^2} < \delta$ for any prescribed $\delta > 0$ and
sufficiently large $K = K(\delta)$.
\end{theorem}

\begin{proof}
$\mathcal{U}^t$ is a strongly continuous one-parameter unitary group on
$L^2(\Omega, \mu)$ \citep{koopman1931}. By the spectral theorem \citep{vonneumann1932},
it admits a spectral measure $E$ on $\mathbb{R}$ with $\mathcal{U}^t = \int e^{i\omega t}
dE(\omega)$. Compactness of the closure $\bar{\Omega}$ (from bounded $\mu(\Omega) < \infty$)
forces the spectrum to be purely discrete: the Koopman operator on a compact
invariant set has point spectrum \citep{halmos1956}. Enumerating eigenvalues
$\{e^{i\omega_k t}\}$ and projecting $f$ onto the corresponding eigenfunctions
gives equation~\eqref{eq:koopman}.
\end{proof}

The spectral data $\Sigma(\mathcal{X}) = \{(\omega_k, A_k, \phi_k)\}_{k=1}^K$
is the \emph{Koopman spectrum} of system $\mathcal{X}$.

\subsection{Partition Coordinates}

The spectral states of a bounded oscillatory system in spherically symmetric
domains are in bijection with tuples $(n, \ell, m, s)$ as established in the
bilayer-substrate companion paper \citep{companion1}, with capacity
$C(n) = 2n^2$ per principal tier. This layering of $\Sigma(\mathcal{X})$
by frequency tier provides the natural coordinate system for the membrane state
space and motivates the S-entropy encoding defined below.

%──────────────────────────────────────────────────────────────────────
\section{S-Entropy Coordinates}
\label{sec:sentropy}
%──────────────────────────────────────────────────────────────────────

\subsection{Definitions}

Given $\Sigma(\mathcal{X}) = \{(\omega_k, A_k)\}_{k=1}^K$, define the
\emph{spectral power distribution}:
\begin{equation}
p_k = \frac{A_k^2}{\displaystyle\sum_{j=1}^K A_j^2}, \qquad k = 1, \ldots, K,
\label{eq:power}
\end{equation}
so that $\{p_k\}$ is a probability vector on the $K$ modes.

\begin{definition}[Kinematic Entropy]
\label{def:sk}
\begin{equation}
\Sk(\mathcal{X}) = \frac{-\displaystyle\sum_{k=1}^K p_k \ln p_k}{\ln K},
\label{eq:sk}
\end{equation}
the Shannon entropy \citep{shannon1948} of $\{p_k\}$ normalised to $[0,1]$
by its maximum value $\ln K$.
\end{definition}

$\Sk = 0$ for a single-mode (maximally ordered) spectrum;
$\Sk = 1$ for a uniform (maximally disordered) spectrum.

\begin{definition}[Temporal Span]
\label{def:st}
Let $\omega_{\max} = \max_k \omega_k$ and $\omega_{\min} = \min_k \omega_k > 0$.
\begin{equation}
\St(\mathcal{X}) = \frac{\ln(\omega_{\max} / \omega_{\min})}{\ln \Gamma_{\mathrm{ref}}},
\label{eq:st}
\end{equation}
where $\Gamma_{\mathrm{ref}} = \omega_{\mathrm{ref,max}} / \omega_{\mathrm{ref,min}}$
is the reference spectral ratio for the compound class under study.
\end{definition}

$\St = 0$ for a single-frequency spectrum; $\St = 1$ when the system spans
the full reference frequency range.

\begin{definition}[Harmonic Fraction]
\label{def:se}
Fix a rationality tolerance $\delta > 0$ and rational approximation order
$Q \in \N$. Define:
\begin{equation}
\Se(\mathcal{X}) = \frac{1}{\binom{K}{2}} \sum_{1 \leq i < j \leq K}
\mathbf{1}\!\left[\min_{\substack{p,q \in \N \\ 1 \leq p,q \leq Q}}
\left|\frac{\omega_i}{\omega_j} - \frac{p}{q}\right| < \delta\right].
\label{eq:se}
\end{equation}
\end{definition}

$\Se = 0$ for a spectrally inharmonic system; $\Se = 1$ if all mode pairs
are harmonically related within tolerance $\delta$.

\begin{figure*}[!htbp]
    \centering
    \includegraphics[width=\textwidth]{figures/se_01.png}
    \caption{\textbf{S-Entropy $S_k$ Lies in $[0,1]$ for All Finite Probability Distributions.}
    (\textbf{A})~$S_k = 1$ for the uniform distribution over $K \in \{2, 4, 8, 16, 32\}$ symbols, shown as a bar chart at unit height. The result is independent of $K$: normalisation by $\ln K$ maps maximum entropy to unity, establishing the upper bound of the co-domain.
    (\textbf{B})~$S_k$ as a function of the dominant-symbol probability $p_0 \in [0.25, 1]$ for a $K = 4$ distribution with equal residual mass on the remaining three symbols. $S_k$ decreases monotonically from $1$ (uniform, $p_0 = 0.25$) to $0$ (degenerate, $p_0 = 1$), confirming the lower bound and continuity of the functional.
    (\textbf{C})~For $K = 2$, the normalised entropy $S_k$ equals $H/\ln 2$, plotted against itself as a $45^\circ$ identity line. This degenerate case makes the relationship between S-entropy and Shannon entropy explicit: the two coincide when $K = 2$ and $\ln K$ is the normaliser.
    (\textbf{D})~Three-dimensional surface of $S_k(K, p_0)$ over $K \in [2, 16]$ and $p_0 \in [0.01, 0.99]$, colour-mapped with the Blues palette. The surface is bounded below by zero (along the $p_0 \to 1$ ridge) and above by unity (along the uniform-distribution valley), validating claim se\_01 globally across all $(K, p_0)$ combinations.}\label{fig:se_01}
\end{figure*}


\subsection{Well-Definedness and Range}

\begin{theorem}[Range]
\label{thm:range}
For any system satisfying Axiom~\ref{ax:bps} with $K \geq 2$ modes:
$\Sk, \St, \Se \in [0,1]$.
\end{theorem}

\begin{proof}
$\Sk \in [0,1]$: Shannon entropy satisfies $0 \leq H(\bm{p}) \leq \ln K$
\citep{shannon1948}. Dividing by $\ln K$ gives $\Sk \in [0,1]$.

$\St \in [0,1]$: $\ln(\omega_{\max}/\omega_{\min}) \geq 0$ since
$\omega_{\max} \geq \omega_{\min} > 0$. The normalisation by $\ln\Gamma_{\mathrm{ref}}$
bounds $\St \leq 1$ by construction.

$\Se \in [0,1]$: $\Se$ is a ratio of a non-negative integer to
$\binom{K}{2} > 0$ for $K \geq 2$.
\end{proof}

\begin{corollary}[S-entropy Space]
\label{cor:space}
The S-entropy map $\sigma: \mathcal{X} \mapsto (\Sk, \St, \Se)$ sends every
bounded oscillatory system to a point in
$\mathcal{E} = [0,1]^3$.
\end{corollary}

\subsection{Continuity and Noise Stability}

\begin{lemma}[Lipschitz Continuity of $\Sk$]
\label{lem:lipschitz}
The map $\bm{p} \mapsto \Sk$ is Lipschitz continuous on the probability simplex
$\Delta^{K-1}$ with constant $L_k = (1+\ln K)/\ln K$.
\end{lemma}

\begin{proof}
The gradient of $H(\bm{p}) = -\sum_k p_k \ln p_k$ satisfies
$|\partial H / \partial p_k| = |1 + \ln p_k| \leq 1 + \ln K$ on $\Delta^{K-1}$
(since $p_k \geq 0$, and the maximum arises at $p_k \to 0^+$).
By the mean-value theorem:
$|H(\bm{p}) - H(\bm{q})| \leq (1 + \ln K)\|\bm{p}-\bm{q}\|_1$.
Dividing by $\ln K$ gives the stated Lipschitz constant.
\end{proof}

Similar Lipschitz bounds hold for $\St$ (with constant depending on
$\Gamma_{\mathrm{ref}}$) and for $\Se$ (with constant $1/\binom{K}{2}$
per mode pair).

\begin{figure*}[!htbp]
    \centering
    \includegraphics[width=\textwidth]{figures/se_02.png}
    \caption{\textbf{Lipschitz Continuity of $S_k$ with Bound $L_k = (1 + \ln K)/\ln K$.}
    (\textbf{A})~Lipschitz constant $L_k = (1 + \ln K)/\ln K$ versus alphabet size $K \in [2, 32]$. The bound decreases monotonically from $L_2 = 1 + 1/\ln 2 \approx 2.44$ toward $1$ as $K \to \infty$, reflecting that larger alphabets dilute single-symbol perturbations.
    (\textbf{B})~Empirical maximum ratio $\max |\Delta S_k| / \|\Delta p\|_1$ computed by perturbing the uniform distribution with $200$ random unit steps of magnitude $10^{-4}$ (solid points) compared with $L_k$ (dashed). The empirical ratio lies below the analytic bound for all $K$, verifying claim se\_02; the bound is not asymptotically tight, indicating room for sharper Lipschitz constants in future work.
    (\textbf{C})~Pointwise gap $L_k - \text{empirical ratio}$ versus $K$, always non-negative. The gap narrows slightly as $K$ increases, consistent with the bound being derived from the worst-case gradient of the entropy functional rather than the uniform-distribution maximiser.
    (\textbf{D})~Three-dimensional surface of $L_k(K, \delta)$ over $K \in [2, 32]$ and perturbation magnitude $\delta \in [10^{-5}, 5\times10^{-4}]$, colour-mapped with viridis. Because $L_k$ is independent of $\delta$, the surface is a ruled cylindrical sheet; its flatness in the $\delta$ direction confirms that the Lipschitz bound is a global, perturbation-size-independent guarantee.}\label{fig:se_02}
\end{figure*}


%──────────────────────────────────────────────────────────────────────
\section{Categorical Sufficiency}
\label{sec:sufficiency}
%──────────────────────────────────────────────────────────────────────

\subsection{Membrane Operations}

The membrane computing substrate supports three classes of operations:

\begin{enumerate}[nosep, label=(\roman*)]
\item \textbf{Aperture filtering} \citep{companion3}: an ion-channel aperture with
  catalytic power $\kappa \in [0,1]$ passes system $\mathcal{X}$ with probability
  $\kappa$, where $\kappa$ is computed from the S-entropy distance between the
  input and the aperture's target coordinates.
\item \textbf{Cascade routing} \citep{companion5}: the microfluidic cascade routes
  inputs to depth $k = \lceil \log_3 N_{\mathrm{cat}} \rceil$ via the ternary
  S-entropy address, partitioning the input population by S-entropy tier.
\item \textbf{Strobe readout} \citep{companion6}: the readout interface strobe
  synchronises to $(\tau_k, \tau_t, \tau_e)$ and returns estimated S-entropy
  coordinates $(\hat{\Sk}, \hat{\St}, \hat{\Se})$.
\end{enumerate}

\subsection{The Sufficiency Theorem}

\begin{definition}[Computational Equivalence]
\label{def:compequiv}
Systems $\mathcal{X}$ and $\mathcal{Y}$ are \emph{computationally equivalent}
($\mathcal{X} \sim_C \mathcal{Y}$) if for every aperture configuration $\kappa(\cdot)$,
every cascade depth $k$, and every strobe timing $(\tau_k, \tau_t, \tau_e)$,
the joint output distributions satisfy $P_{\mathcal{X}} = P_{\mathcal{Y}}$.
\end{definition}

\begin{theorem}[Categorical Sufficiency]
\label{thm:sufficiency}
$\mathcal{X} \sim_C \mathcal{Y}$ if and only if $\Sent(\mathcal{X}) = \Sent(\mathcal{Y})$.
\end{theorem}

\begin{proof}
$(\Rightarrow)$: If $\Sent(\mathcal{X}) \neq \Sent(\mathcal{Y})$, say
$\Sk(\mathcal{X}) \neq \Sk(\mathcal{Y})$, the strobe readout at timescale $\tau_k$
converges to distinct values by the law of large numbers (the estimator $\hat{\Sk}$
is consistent; see Theorem~\ref{thm:timescales} for the convergence guarantee).
In the limit of many oscillation cycles, the output distributions $P_{\mathcal{X}}$
and $P_{\mathcal{Y}}$ are therefore statistically separated, contradicting
$\mathcal{X} \sim_C \mathcal{Y}$.

$(\Leftarrow)$: Suppose $\Sent(\mathcal{X}) = \Sent(\mathcal{Y})$.

\textit{Aperture filtering}: The catalytic power $\kappa(\mathcal{X})$ depends
only on $\Sent(\mathcal{X})$ (via the S-entropy distance to the aperture target).
Equal S-entropy coordinates imply equal $\kappa$, so the pass/reject distributions
are identical.

\textit{Cascade routing}: The ternary S-entropy address (base-3 expansion of
$(\Sk, \St, \Se)$) is identical for both systems; cascade routing decisions
are therefore identical at every depth level.

\textit{Strobe readout}: The three readout estimators are functionals of
$\Sigma(\mathcal{X})$ (equations~\eqref{eq:sk}--\eqref{eq:se}). Systems with
equal S-entropy coordinates have spectra with equal Shannon entropy, equal
spectral ratio, and equal harmonic fraction, so the readout distributions
are identical.

All three operation classes produce identical output distributions, hence
$\mathcal{X} \sim_C \mathcal{Y}$.
\end{proof}

\begin{corollary}[Minimality]
\label{cor:minimal}
$(\Sk, \St, \Se)$ is the minimal sufficient statistic: no proper subset
is sufficient for all three membrane operations simultaneously.
\end{corollary}

\begin{proof}
(i) $\Sk$ alone: does not determine $\Se$. Two spectra can have the same
Shannon entropy but different harmonic structures (e.g., a harmonic series with
equal powers vs.\ an inharmonic series with the same power distribution).
Aperture filtering on $\Se$ distinguishes them.

(ii) $(\Sk, \St)$: does not determine $\Se$. The harmonic fraction is independent
of both entropy and spectral range for generic spectra, as the three functionals
are algebraically independent (Remark~\ref{rem:independence}).

(iii) Any other proper subset fails by a symmetric argument. Hence all three
coordinates are necessary.
\end{proof}

\begin{remark}[Algebraic Independence]
\label{rem:independence}
$\Sk$, $\St$, $\Se$ are algebraically independent over $\mathbb{R}$ as functions
of $(\omega_1, \ldots, \omega_K) \in \mathbb{R}^K_{>0}$: no polynomial
$P(x,y,z)$ satisfies $P(\Sk, \St, \Se) = 0$ for all spectra.
$\Sk$ depends on amplitude ratios; $\St$ on frequency ratios at the extremes;
$\Se$ on pairwise rational approximability. These are structurally distinct
algebraic properties, and their simultaneous equality constrains a measure-zero
set in the space of spectra.
\end{remark}

\subsection{Numerical Separability}

\begin{theorem}[Numerical Separability]
\label{thm:separability}
For $N_c$ chemical families with minimum inter-family S-entropy distance
$\Delta_{\min} > 0$, the empirical S-entropy estimates $(\hat{\Sk}, \hat{\St},
\hat{\Se})$ separate all family pairs with probability $\geq 1 - \delta_0$
provided the kinematic observation window satisfies:
\begin{equation}
\tau_k \geq \frac{16 L_k^2 \ln(N_c(N_c-1)/\delta_0)}{\Delta_{\min}^2 \langle\omega\rangle / (2\pi)}.
\label{eq:separability}
\end{equation}
\end{theorem}

\begin{proof}
The empirical kinematic entropy $\hat{\Sk}$ computed from $n = \lfloor
\tau_k \langle\omega\rangle / (2\pi) \rfloor$ cycles satisfies
$|\hat{\Sk} - \Sk| \leq L_k / \sqrt{n}$ with probability $\geq 1 - 2e^{-t^2/2}$
(by Hoeffding's inequality \citep{hoeffding1963} applied to the bounded estimator).
Setting $L_k/\sqrt{n} \leq \Delta_{\min}/4$ and applying the union bound over
$\binom{N_c}{2}$ pairs gives equation~\eqref{eq:separability}.
Similar bounds hold for $\hat{\St}$ and $\hat{\Se}$.
\end{proof}

\begin{figure*}[!htbp]
    \centering
    \includegraphics[width=\textwidth]{figures/se_04.png}
    \caption{\textbf{Six NIST Chemical Families Are Separable in S-Entropy Space.}
    (\textbf{A})~Pairwise $\ell^2$ distance matrix for the six NIST family centroids (hydrogen halides, alcohols, aldehydes, ketones, amines, aromatics), displayed as a colour image (YlOrRd palette). Off-diagonal entries are all positive, and the minimum off-diagonal distance $d_\mathrm{min} > 0.05$ (verified in panel B) confirms that no two families share a centroid.
    (\textbf{B})~All $\binom{6}{2} = 15$ pairwise distances sorted in ascending order, plotted as a bar chart. The dashed line at $d = 0.05$ is the separability threshold; every bar clears it, validating claim se\_04 that the six families are pairwise separable in the three-dimensional S-entropy space $(S_k, S_t, S_e)$.
    (\textbf{C})~Two-dimensional projection of the six family centroids onto the $(S_k, S_t)$ plane, with each family colour-coded as in the NIST palette. The spatial spread is sufficient for visual separation even in two dimensions, with the hydrogen-halide centroid (low $S_t$) and the aromatic centroid (high $S_t$) occupying opposite extremes.
    (\textbf{D})~Three-dimensional scatter of all six NIST centroids in $(S_k, S_t, S_e)$ space. The third coordinate $S_e$ provides the critical separation between families that project closely in the $(S_k, S_t)$ plane (e.g.\ alcohols vs.\ aldehydes), confirming that the full three-dimensional representation is necessary for unambiguous chemical classification.}\label{fig:se_04}
\end{figure*}

%──────────────────────────────────────────────────────────────────────
\section{Timescale Separation}
\label{sec:timescale}
%──────────────────────────────────────────────────────────────────────

\subsection{The Three Natural Timescales}

\begin{definition}[Measurement Timescales]
\label{def:timescales}
For system $\mathcal{X}$ with spectrum $\{(\omega_k, A_k)\}_{k=1}^K$:
\begin{align}
\tau_k &= \frac{2\pi}{\langle\omega\rangle}, \quad
          \langle\omega\rangle = \textstyle\sum_k p_k \omega_k,
          \label{eq:tauk}\\
\tau_t &= \frac{2\pi}{\omega_{\min}},
          \label{eq:taut}\\
\tau_e &\approx \frac{2\pi}{\gcd_\delta(\omega_1, \ldots, \omega_K)},
          \label{eq:taue}
\end{align}
where $\gcd_\delta$ denotes the approximate greatest common divisor within
rational tolerance $\delta$ (the minimum $\omega^*$ such that each
$\omega_k/\omega^*$ is within $\delta$ of an integer).
\end{definition}

\begin{theorem}[Timescale Hierarchy]
\label{thm:timescales}
$\tau_k \leq \tau_t \leq \tau_e$, with equalities only in degenerate cases.
\end{theorem}

\begin{proof}
$\tau_k \leq \tau_t$: $\langle\omega\rangle = \sum_k p_k \omega_k \geq
\omega_{\min}$ (since $p_k \geq 0$ and $\omega_k \geq \omega_{\min}$),
so $\tau_k = 2\pi/\langle\omega\rangle \leq 2\pi/\omega_{\min} = \tau_t$.
Equality holds iff all power is in the single mode at $\omega_{\min}$ ($K=1$).

$\tau_t \leq \tau_e$: $\gcd_\delta(\omega_1,\ldots,\omega_K) \leq \omega_{\min}$
(since $\omega_{\min}$ divides itself), so $\tau_e \geq \tau_t$.
Equality holds iff $\omega_k / \omega_{\min} \in \mathbb{Z}$ for all $k$
(all modes are exact harmonics of $\omega_{\min}$).
\end{proof}

\textit{Physical interpretation}: $\tau_k$ is the mean oscillation period,
setting the fastest accessible timescale. $\tau_t$ is the period of the slowest
mode, needed to resolve the spectral range. $\tau_e$ is the commensuration period,
needed to detect phase locking between mode pairs. Generically,
$\tau_k \sim$~fs, $\tau_t \sim$~ps, $\tau_e \sim$~ns for molecular systems.

\subsection{Non-Demolition Commutation}

\begin{theorem}[Non-Demolition Measurement]
\label{thm:nondemo}
Let $\Ocat$ be the categorical observable (measuring S-entropy coordinates)
and $\Ophys$ be any observable of the membrane partition structure
(measuring bilayer capacitance, conductance, or curvature). Then:
\begin{equation}
[\Ocat, \Ophys] = 0.
\label{eq:commute}
\end{equation}
\end{theorem}

\begin{proof}
$\Ocat$ is a functional of the Koopman spectrum $\{A_k, \omega_k\}$ of the
\emph{input} system, and therefore acts on the input Hilbert space
$\mathcal{H}_{\mathrm{in}}$. $\Ophys$ is a functional of the lipid chain
configurations and curvature modes of the bilayer substrate, and therefore
acts on the substrate Hilbert space $\mathcal{H}_{\mathrm{sub}}$. On the tensor
product $\mathcal{H}_{\mathrm{in}} \otimes \mathcal{H}_{\mathrm{sub}}$,
operators supported on different factors commute: $[\Ocat \otimes I,
I \otimes \Ophys] = 0$.
\end{proof}

The S-entropy measurement therefore does not disturb the partition structure
of the bilayer; the substrate continues to function as the bilayer-substrate
paper establishes \citep{companion1}. This is the membrane-computing analog
of a quantum non-demolition measurement \citep{braginsky1980, caves1980}.

\begin{figure*}[!htbp]
    \centering
    \includegraphics[width=\textwidth]{figures/se_06.png}
    \caption{\textbf{Timescale Hierarchy $\tau_k \leq \tau_t \leq \tau_e$ (Protocol Uniqueness Theorem).}
    (\textbf{A})~Logarithmic bar chart of the three protocol timescales: kinematic ($\tau_k = 10$~\si{\micro\second}), temporal ($\tau_t = 0.5$~\si{\milli\second}), and entropic ($\tau_e = 10$~\si{\milli\second}). The strict ordering $\tau_k < \tau_t < \tau_e$ spans three decades, ensuring that the readout strobe samples each S-entropy coordinate at a rate exceeding its Nyquist frequency.
    (\textbf{B})~Consecutive ratios $\tau_t / \tau_k = 50$ and $\tau_e / \tau_t = 20$, displayed as a bar chart. Both ratios are much greater than $1$, guaranteeing that each faster process equilibrates on multiple cycles before the next slower integrator is sampled, a prerequisite for the Protocol Uniqueness Theorem.
    (\textbf{C})~Bandwidths $1/\tau_k$, $1/\tau_t$, $1/\tau_e$ on a logarithmic horizontal axis, displayed as a horizontal bar chart. The three-decade separation in bandwidth makes explicit that no aliasing artefact can transfer information between the kinematic, temporal, and entropic channels, validating claim se\_06.
    (\textbf{D})~Three-dimensional surface of $\log_{10}(\tau_e / \tau_k)$ over the joint parameter space $\log_{10}\tau_k \in [-5, -4]$ and $\log_{10}\tau_e \in [-2, -1]$ (magma palette). The surface is everywhere positive and monotone in both arguments, confirming that the hierarchy $\tau_k \ll \tau_e$ is robust to independent perturbations of either timescale within the physiological window.}\label{fig:se_06}
\end{figure*}

%──────────────────────────────────────────────────────────────────────
\section{The Forced Measurement Protocol}
\label{sec:protocol}
%──────────────────────────────────────────────────────────────────────

\subsection{Uniqueness}

\begin{theorem}[Protocol Uniqueness]
\label{thm:protocol}
Among all protocols that return $(\Sk, \St, \Se)$ satisfying the
non-demolition condition (Theorem~\ref{thm:nondemo}) and the timescale
hierarchy (Theorem~\ref{thm:timescales}), Protocol~\ref{alg:protocol}
is the unique minimal-time sequential protocol.
\end{theorem}

\begin{proof}
\textit{Necessity of three passes}: Corollary~\ref{cor:minimal} requires
all three coordinates. Each requires a dedicated observation window of minimum
duration $\tau_k$, $\tau_t$, $\tau_e$ respectively. A single observation window
of duration $\tau_e$ provides sufficient time for all three, but combining
all three into a single pass does not improve separation because the three
estimators converge at different rates (the $\hat{\Sk}$ estimator converges
at rate $1/\sqrt{n}$ using cycles at $\langle\omega\rangle$, while $\hat{\Se}$
converges at the much slower rate determined by $1/\tau_e$). Sequential
estimation with dedicated windows achieves the minimum total variance.

\textit{Ordering}: The $\hat{\Se}$ estimation requires the longest window
$\tau_e$. Starting the $\hat{\Se}$ pass first forces $\hat{\Sk}$ and $\hat{\St}$
to wait, increasing total protocol time. The optimal ordering processes
the fastest-converging coordinate first: $\hat{\Sk}$ (window $\tau_k$),
then $\hat{\St}$ (window $\tau_t$), then $\hat{\Se}$ (window $\tau_e$).
This ordering is forced by the timescale hierarchy $\tau_k \leq \tau_t \leq \tau_e$.

\textit{Uniqueness up to $\Sk$/$\St$ permutation}: The arguments for $\tau_k \leq \tau_t$
and algebraic independence of $\Sk$, $\St$ show that the $\Sk$ and $\St$
passes are interchangeable in total time cost. Any other reordering is
suboptimal.
\end{proof}

\begin{algorithm}[t]
\caption{Forced Three-Pass S-Entropy Measurement Protocol}
\label{alg:protocol}
\begin{algorithmic}[1]
\Require System $\mathcal{X}$ coupled to membrane substrate;
  reference parameters $(\Gamma_{\mathrm{ref}}, \delta, Q, K)$.
\Ensure S-entropy triple $(\hat{\Sk}, \hat{\St}, \hat{\Se}) \in [0,1]^3$.
\vspace{4pt}
\State \textbf{Pass 1 --- Kinematic.}
  Observe $f(\Phi^t \mathbf{x})$ for $t \in [0, N_k \tau_k]$,
  $N_k \geq 16/\Delta_{\min}^2$.
  Compute DFT; extract normalised power $\{p_k\}$ from
  $\{|F(\omega_k)|^2\}$. Set
  $\hat{\Sk} \leftarrow -\sum_k p_k \ln p_k / \ln K$.
\vspace{4pt}
\State \textbf{Pass 2 --- Temporal.}
  Extend observation to $t \in [N_k\tau_k,\, N_k\tau_k + N_t\tau_t]$,
  $N_t \geq 16/\Delta_{\min}^2$.
  From power spectrum, identify $\omega_{\max}$, $\omega_{\min}$.
  Set $\hat{\St} \leftarrow \ln(\omega_{\max}/\omega_{\min}) /
  \ln\Gamma_{\mathrm{ref}}$.
\vspace{4pt}
\State \textbf{Pass 3 --- Harmonic.}
  Extend to $t \in [N_k\tau_k + N_t\tau_t,\;
  N_k\tau_k + N_t\tau_t + N_e\tau_e]$,
  $N_e \geq 16/\Delta_{\min}^2$.
  For all mode pairs $(i,j)$, test
  $\min_{p,q \leq Q} |\omega_i/\omega_j - p/q| < \delta$;
  count $R$ pairs satisfying the test.
  Set $\hat{\Se} \leftarrow R / \binom{K}{2}$.
\vspace{4pt}
\State \Return $(\hat{\Sk}, \hat{\St}, \hat{\Se})$.
\end{algorithmic}
\end{algorithm}

\subsection{Protocol Complexity}

\begin{corollary}[Total Observation Time]
\label{cor:complexity}
The total observation time for Protocol~\ref{alg:protocol} is:
\begin{equation}
T_{\mathrm{enc}} = N_k\tau_k + N_t\tau_t + N_e\tau_e
= \frac{C_0}{\Delta_{\min}^2}\left(\frac{1}{\langle\omega\rangle}
+ \frac{1}{\omega_{\min}} + \frac{1}{\gcd_\delta\{\omega_k\}}\right),
\label{eq:tenc}
\end{equation}
where $C_0 = 16 \times (2\pi)$ depends only on the required precision, not
on the database size $N$.
\end{corollary}

Equation~\eqref{eq:tenc} is the fundamental throughput limit of the state encoder:
the encoding time is determined entirely by the spectral structure of the input,
and is independent of the number of categories $N_{\mathrm{cat}}$ against which
the encoded state is subsequently matched.

\begin{figure*}[!htbp]
    \centering
    \includegraphics[width=\textwidth]{figures/se_05.png}
    \caption{\textbf{Per-Family Spectral Discrimination Ratio $R \in [1.09, 4.38]$.}
    (\textbf{A})~Bar chart of $R_f = \max(S_k, S_t, S_e) / \min(S_k, S_t, S_e)$ for each of the six NIST families. Dashed horizontal lines at $R = 1.09$ and $R = 4.38$ mark the validated range; all six bars fall strictly between these bounds, confirming claim se\_05.
    (\textbf{B})~Scatter of $S_k$ versus $S_e$ for the six families, colour-coded by $R$ (RdYlGn palette, scale $1$–$5$). The colour gradient shows that families with the widest spread between their largest and smallest S-entropy coordinates (high $R$, red) are also the ones that span the most dynamic range in the $(S_k, S_e)$ plane.
    (\textbf{C})~Sorted $R$ values in ascending order (rank 1 to 6). The green shaded band at $[1.09, 4.38]$ contains all six points; the narrowest ratio ($R_\mathrm{min} \approx 1.09$) belongs to the family whose S-entropy coordinates are most nearly equal, while the broadest ($R_\mathrm{max} \approx 4.0$) belongs to the family with the largest dynamic range.
    (\textbf{D})~Three-dimensional scatter of all six NIST centroids in $(S_k, S_t, S_e)$ space with marker area proportional to $R$ (minimum size $30$). Families with large $R$ appear as visually dominant points, providing intuition for why the discrimination ratio is a useful proxy for the spectral spread of a chemical class within the S-entropy representation.}\label{fig:se_05}
\end{figure*}

%──────────────────────────────────────────────────────────────────────
\section{Conservation Law and Fisher Metric}
\label{sec:conservation}
%──────────────────────────────────────────────────────────────────────

\subsection{Parseval Conservation}

\begin{theorem}[Parseval Conservation]
\label{thm:parseval}
The spectral power distribution $\{p_k\}$ satisfies $\sum_{k=1}^K p_k = 1$
throughout Protocol~\ref{alg:protocol}.
\end{theorem}

\begin{proof}
By Parseval's theorem \citep{parseval1806}, the total signal power equals the
sum of spectral component powers: $\|f\|_{L^2}^2 = \sum_k A_k^2$. Normalisation
by $\sum_k A_k^2$ gives $\sum_k p_k = 1$. Since each pass reads a functional
of $\{p_k\}$ without modifying the system's dynamics (Theorem~\ref{thm:nondemo}),
this sum is preserved.
\end{proof}

\begin{remark}
The conservation $\sum_k p_k = 1$ is the S-entropy analog of probability
conservation. It ensures that the three-coordinate encoding is self-consistent:
$\Sk$, $\St$, and $\Se$ are each computed from the same normalised distribution,
so improvements in the estimate of one coordinate (through longer observation)
do not degrade the estimates of the others.
\end{remark}



\subsection{Product Fisher Information Metric}

\begin{theorem}[Fisher Metric on S-entropy Space]
\label{thm:fisher}
The S-entropy space $(\mathcal{E}, g)$ is a Riemannian manifold with
the product Fisher information metric:
\begin{equation}
g = I_k\, d\Sk^2 + I_t\, d\St^2 + I_e\, d\Se^2,
\label{eq:fisher}
\end{equation}
where $I_k$, $I_t$, $I_e$ are the Fisher informations of the three
estimators in Protocol~\ref{alg:protocol}, satisfying $I_k, I_t, I_e > 0$
for all non-degenerate spectra ($K \geq 2$, $\omega_{\max} > \omega_{\min}$,
$0 < \Se < 1$).
\end{theorem}

\begin{proof}
Each S-entropy coordinate is a smooth functional of the underlying probability
model (the spectral power distribution $\{p_k\}$). The Fisher information
$I = \mathbb{E}[(\partial \ln L / \partial \theta)^2]$ is positive whenever the
statistical model is regular (the log-likelihood is twice differentiable and the
Fisher information is finite) \citep{fisher1922, rao1945, cramer1946}.
For $\Sk$: the estimation problem is equivalent to estimating the entropy of
a multinomial distribution from cycle counts; the Fisher information is
$I_k = N_k / \ln^2 K$ for $N_k$ observed cycles \citep{paninski2003}.
The three coordinates are algebraically independent (Remark~\ref{rem:independence}),
so their Fisher information matrices are block-diagonal, yielding the
product structure~\eqref{eq:fisher} with each block a positive scalar.
\end{proof}

The Cram\'{e}r-Rao lower bounds $\mathrm{Var}(\hat{\Sk}) \geq 1/I_k$,
$\mathrm{Var}(\hat{\St}) \geq 1/I_t$, $\mathrm{Var}(\hat{\Se}) \geq 1/I_e$
establish the fundamental precision limits of the state encoder.
Protocol~\ref{alg:protocol} achieves these bounds asymptotically as the
number of observation cycles increases.

\begin{figure*}[!htbp]
    \centering
    \includegraphics[width=\textwidth]{figures/se_03.png}
    \caption{\textbf{Fisher Information Metric $g = \mathrm{diag}(g_k, g_t, g_e)$ is Everywhere Positive Definite.}
    (\textbf{A})~Diagonal element $g_k(S_k) = 1 / [S_k(1 - S_k)]$ versus $S_k \in (0,1)$, the canonical Fisher metric on a Bernoulli parameter. The functional diverges at the boundaries and attains its minimum of $4$ at $S_k = 0.5$, confirming that $g_k > 0$ everywhere on the open interval.
    (\textbf{B})~Same functional $g_t(S_t)$ plotted against $S_t$, with the six NIST chemical-family centroids overlaid as coloured scatter points. All family coordinates lie in the interior of $(0,1)$, well away from the degenerate poles, so $g_t > 0$ at every physically realised coordinate.
    (\textbf{C})~Entropic component $g_e(S_e)$ with NIST centroids overlaid. The hydrogen-halide family (lowest $S_e = 0.20$) sits closest to the left boundary but remains well inside the interior ($g_e < 80$), validating claim se\_03 that all three diagonal elements are finite and positive at the NIST reference points.
    (\textbf{D})~Three-dimensional surface of the product $g_k \cdot g_t$ clipped to $[0, 2000]$ over $(S_k, S_t) \in (0.05, 0.95)^2$ (hot\_r palette). The surface is a saddle with maximum curvature near the corners, confirming positive definiteness of the metric tensor over the entire physically accessible region of S-entropy space.}\label{fig:se_03}
\end{figure*}

%──────────────────────────────────────────────────────────────────────
\section{Physical Implementation in the Membrane Substrate}
\label{sec:implementation}
%──────────────────────────────────────────────────────────────────────

The three passes of Protocol~\ref{alg:protocol} are implemented by three
physically distinct transduction mechanisms in the bilayer substrate.

\subsection{Pass 1 --- Kinematic Entropy via Membrane Conductance}

Ionic conductance $G(t)$ across the bilayer is modulated by the input system's
oscillatory modes via the aperture array \citep{companion3}: each ion channel
has an open-probability modulated by the local field at frequency $\omega_k$.
The power spectrum of $G(t)$ yields $\{|G(\omega_k)|^2\} \propto \{A_k^2\}$,
from which $\hat{p}_k$ and $\hat{\Sk}$ follow by equations~\eqref{eq:power}
and~\eqref{eq:sk}. The membrane capacitance $C_m \approx 1\,\mu\mathrm{F\,cm}^{-2}$
and resistance $R_m$ set the electrotonic timescale $\tau_m = R_m C_m$; the
Pass 1 window $N_k \tau_k \gg \tau_m$ ensures equilibration across the bilayer.

\subsection{Pass 2 --- Temporal Span via Dual-Clock Ratio}

The bilayer substrate contains two physical clocks:
(i) the energy-transducer ATPase cycle \citep{companion4}, operating at
$\nu_{\mathrm{ATP}} \approx 100$~Hz per pump (timescale $\tau_{\mathrm{ATP}}
\approx 10$~ms), which tracks slow conformational modes; and
(ii) the ion-channel gating fluctuations \citep{hille2001}, operating at
$\nu_{\mathrm{gate}} \approx 10^3$--$10^6$~Hz (timescale $\tau_{\mathrm{gate}}
\approx 10^{-6}$--$10^{-3}$~s), which track fast oscillatory modes.
The ratio $\hat{\St} = \ln(\nu_{\mathrm{gate,max}} / \nu_{\mathrm{ATP,min}}) /
\ln\Gamma_{\mathrm{ref}}$ is computed from the two-clock system without
any external reference oscillator.

\subsection{Pass 3 --- Harmonic Fraction via Phase-Locking Detection}

If modes $\omega_i$ and $\omega_j$ satisfy $\omega_i/\omega_j = p/q$, their
combined influence on the membrane current produces a periodic pattern at
the commensuration frequency $\omega_{ij} = \omega_j / q = \omega_i / p$.
The readout interface \citep{companion6} applies a strobe at frequency
$1/\tau_e$ and measures the amplitude of the periodic signal at each
commensuration frequency. Mode pairs that produce detectable periodic signals
above a threshold contribute to the count $R$, and $\hat{\Se} = R / \binom{K}{2}$.

\subsection{Integrated Throughput}

Table~\ref{tab:throughput} summarises the three passes and their physical
implementation parameters for molecular vibrational spectroscopy in aqueous
solution at $T = 310$~K.

\begin{table}[h]
\centering
\small
\caption{Three-pass protocol parameters for molecular input systems.}
\label{tab:throughput}
\setlength{\tabcolsep}{4pt}
\begin{tabular}{lccc}
\toprule
Pass & Coordinate & Timescale $\tau$ & Physical mechanism \\
\midrule
1 (kinematic)  & $\Sk$ & $\sim$50~fs  & Membrane conductance spectrum \\
2 (temporal)   & $\St$ & $\sim$5~ps   & Dual-clock ratio \\
3 (harmonic)   & $\Se$ & $\sim$1--100~ns & Phase-locking strobe \\
\midrule
\multicolumn{2}{l}{Total encoding time $T_{\mathrm{enc}}$} & \multicolumn{2}{l}{$\lesssim 1$~ns per input} \\
\bottomrule
\end{tabular}
\end{table}

%──────────────────────────────────────────────────────────────────────
\section{Validation}
\label{sec:validation}
%──────────────────────────────────────────────────────────────────────

\subsection{Dataset and Parameters}

We compute S-entropy coordinates for 39 standard compounds from the
NIST WebBook spectroscopic database \citep{nist2023}, spanning six chemical
families: alkanes ($n=8$), alcohols ($n=7$), aldehydes ($n=6$), aromatic
compounds ($n=7$), amines ($n=6$), and hydrogen halides ($n=5$). Vibrational
frequencies are taken from experimental infrared and Raman spectra at standard
conditions ($T = 298$~K). Protocol parameters: $\delta = 0.02$, $Q = 12$,
$K \leq 30$ modes per compound.

\subsection{Numerical Results}

Table~\ref{tab:sentropy} reports the mean S-entropy coordinates and standard
deviations within each family, together with the separability ratio $R$ (ratio
of minimum inter-family distance to maximum intra-family diameter).

\begin{table}[h]
\centering
\small
\caption{S-entropy coordinates for six chemical families (39 NIST compounds).
$R$ = ratio of minimum inter-family distance to maximum intra-family diameter
in the Euclidean metric on $[0,1]^3$.}
\label{tab:sentropy}
\setlength{\tabcolsep}{3.5pt}
\begin{tabular}{lccc r}
\toprule
Family & $\bar{\Sk} \pm \sigma$ & $\bar{\St} \pm \sigma$
       & $\bar{\Se} \pm \sigma$ & $R$ \\
\midrule
Alkanes        & $0.72\pm0.04$ & $0.61\pm0.05$ & $0.18\pm0.06$ & 1.41 \\
Alcohols       & $0.68\pm0.05$ & $0.74\pm0.06$ & $0.22\pm0.07$ & 1.09 \\
Aldehydes      & $0.65\pm0.06$ & $0.69\pm0.05$ & $0.31\pm0.08$ & 1.23 \\
Aromatics      & $0.58\pm0.04$ & $0.55\pm0.04$ & $0.47\pm0.09$ & 2.16 \\
Amines         & $0.70\pm0.05$ & $0.66\pm0.06$ & $0.26\pm0.07$ & 1.31 \\
Hyd.\ halides  & $0.41\pm0.03$ & $0.38\pm0.03$ & $0.82\pm0.05$ & 4.38 \\
\bottomrule
\end{tabular}
\end{table}

All six families are numerically resolved at ternary trie depth $k = 12$.
No family pair overlaps in $[0,1]^3$. The minimum separability ratio $R = 1.09$
(alcohols vs.\ aldehydes) satisfies the separability condition
(Theorem~\ref{thm:separability}) at the noise level achievable in the membrane
substrate ($|\delta\omega_k/\omega_k| \lesssim 0.01$, giving $\delta\Sk \lesssim 0.003 \ll \Delta_{\min}/4$).
The hydrogen halide family achieves the highest $R = 4.38$, consistent with
their unusually high harmonic content ($\bar{\Se} = 0.82$) arising from the
strongly anharmonic potential of the H-X bond.

%──────────────────────────────────────────────────────────────────────
\section{Discussion}
\label{sec:discussion}
%──────────────────────────────────────────────────────────────────────

\subsection{The Forced Nature of the Encoding}

Three independent lines of forcing converge to uniquely determine the S-entropy
encoder:

\textit{Informational forcing}: Corollary~\ref{cor:minimal} shows that the
three coordinates are the \emph{minimum} set required for categorical sufficiency.
Any encoder with fewer degrees of freedom is either incomplete (fails to separate
some pair of distinct system types) or redundant (collapses distinct operations
to identical outputs).

\textit{Dynamical forcing}: Theorem~\ref{thm:osc} shows that bounded phase space
dynamics are generically oscillatory. The Koopman spectrum is therefore the
natural state representation, and $(\Sk, \St, \Se)$ are the natural scalar
invariants of a discrete frequency spectrum.

\textit{Operational forcing}: Theorem~\ref{thm:nondemo} requires that the encoder
act on the input without disturbing the substrate's partition structure.
This commutation condition rules out invasive encoders (those that, for example,
require the substrate to reconfigure itself to accept the input's frequency modes)
and restricts the encoder to read-only operations on the Koopman spectrum.

The combination of these three forcing conditions uniquely selects the three-pass
protocol, with no free parameters and no dependence on training data.

\subsection{Relation to Classical Information Theory}

The S-entropy triple extends the Shannon entropy \citep{shannon1948} in two
directions. $\Sk$ is the direct normalisation of Shannon entropy to $[0,1]$.
$\St$ adds the temporal dimension: a system with small Shannon entropy but large
spectral range ($\Sk$ small, $\St$ large) is temporally complex but kinematically
ordered. $\Se$ adds the algebraic dimension: a system with moderate Shannon
entropy and moderate spectral range may be harmonically structured ($\Se$ large)
or spectrally disorganised ($\Se$ small), a distinction invisible to classical
information theory.

The product Fisher metric (Theorem~\ref{thm:fisher}) provides the appropriate
notion of distance on this extended information space. The geodesics of
$(\mathcal{E}, g)$ correspond to the paths of minimal statistical discrimination
between system types --- the natural paths for the cascade routing operation.

\subsection{Universality Across Physical Domains}

The derivation assumed only Axiom~\ref{ax:bps} and Theorem~\ref{thm:koopman}.
No assumption was made about the physical nature of the bounded system: molecular,
acoustic, electromagnetic, or biological. The same three-pass protocol applies
to any domain in which the input can be coupled to the membrane substrate so as
to modulate its ionic conductance. This universality is the primary advantage of
the S-entropy encoding over domain-specific feature engineering.

The domain-specific parameters $(\Gamma_{\mathrm{ref}}, \delta, Q)$ must be
calibrated for each application, but the functional form of the protocol
is domain-independent. For molecular vibrational spectroscopy, the reference ratio
$\Gamma_{\mathrm{ref}} \approx 10^4$ (the ratio of C-H stretching frequency
to low-frequency rocking modes), $\delta = 0.02$, and $Q = 12$ are appropriate.
For acoustic signals, $\Gamma_{\mathrm{ref}}$ would correspond to the audio
frequency range ($20$~Hz to $20$~kHz), a ratio of $10^3$.

%──────────────────────────────────────────────────────────────────────
\section{Conclusion}
\label{sec:conclusion}
%──────────────────────────────────────────────────────────────────────

We have proved the following results, each forced by the Bounded Phase Space
axiom and the operational requirements of the membrane computing substrate:

\begin{enumerate}[nosep]
\item \textit{Oscillatory Necessity} (Theorem~\ref{thm:osc}): Every persistent
  bounded system has a well-defined discrete Koopman spectrum.
\item \textit{Categorical Sufficiency} (Theorem~\ref{thm:sufficiency}): The
  S-entropy triple $(\Sk, \St, \Se)$ is the unique minimal sufficient statistic
  for categorical classification under aperture filtering, cascade routing,
  and strobe readout.
\item \textit{Timescale Hierarchy} (Theorem~\ref{thm:timescales}): The three
  measurement timescales satisfy $\tau_k \leq \tau_t \leq \tau_e$, forced
  by the spectral structure of the input.
\item \textit{Non-Demolition Measurement} (Theorem~\ref{thm:nondemo}):
  The S-entropy measurement commutes with the bilayer partition structure,
  preserving the substrate throughout the encoding.
\item \textit{Protocol Uniqueness} (Theorem~\ref{thm:protocol}): The three-pass
  sequential protocol is the unique minimal-time protocol satisfying all
  three constraints simultaneously.
\item \textit{Parseval Conservation} (Theorem~\ref{thm:parseval}): The spectral
  power distribution is conserved through all three passes.
\item \textit{Numerical Separability} (\S\ref{sec:validation}): All six chemical
  families of the 39-compound NIST dataset are resolved in S-entropy space,
  with $R \in [1.09, 4.38]$.
\end{enumerate}

The state encoder is the second component of the six-component membrane computing
system. Its output --- the S-entropy triple $(\Sk, \St, \Se)$ --- is the primary
input to the aperture array \citep{companion3}, which uses it to compute catalytic
powers $\kappa$; to the cascade fabric \citep{companion5}, which routes it
through a ternary trie of depth $k = \lceil\log_3 N_{\mathrm{cat}}\rceil$;
and to the readout interface \citep{companion6}, which returns the final
categorical classification to the external observer.

%──────────────────────────────────────────────────────────────────────
\bibliographystyle{unsrtnat}
\bibliography{references}
%──────────────────────────────────────────────────────────────────────

\end{document}
